\documentclass[letterpaper]{article}
\PassOptionsToPackage{unicode=true}{hyperref} % options for packages loaded elsewhere
\PassOptionsToPackage{hyphens}{url}
\usepackage{graphicx} % Required for inserting images


\def\tightlist{}


\graphicspath{ {../images/} }
\usepackage[margin=1in]{geometry}
\usepackage{tikz}
\usepackage[dvipsnames]{xcolor}
\usetikzlibrary{calc}




\usepackage[utf8]{inputenc}
\usepackage[T1]{fontenc}
\usepackage{lmodern}
\usepackage[english]{babel}

\usepackage{longtable}
\usepackage{booktabs}

% Useful packages
\usepackage{hyperref}
\usepackage{amsmath, amssymb}
\usepackage{fancyhdr}
\usepackage{setspace}
\usepackage{tocloft}
\usepackage{titlesec}


\usepackage{unicode-math}



% Code highlighting (Pandoc uses listings or minted, but listings is safer without shell escape)
\usepackage{listings}
\lstset{
  basicstyle=\ttfamily\small,
  breaklines=true,
  frame=single,
  backgroundcolor=\color[gray]{0.95}
}



%Modifiable Commands; defaults set to none
\newcommand{\titleDOC}{}
\newcommand{\authorDOC}{Daniel Sabalakov}
\newcommand{\dateDOC}{\today}
\newcommand{\classDOC}{}
\newcommand{\emailDOC}{danielsabalakov@gmail.com}
% DEFAULT QUOTE IS BY ISSAC NEWTON. IF YOU DON'T WANT IT, DO: quote: @none
\newcommand{\quoteDOC}{``\textit{Truth is ever to be found in the simplicity, and not in the multiplicity and confusion of things}'' - Issac Newton}

% Header/Footer
\pagestyle{fancy}
\fancyhf{}
\rhead{\thepage}
\lhead{\leftmark}





% %Title Arguments
% \title{\TITLEOFDOCUMENT}
% \author{\AUTHOROFDOCUMENT}
% \date{\today}


%
% ANYTHING BELOW HERE CAN GET PROCEDURALLY GENERATED
%
% BEGINNING OF DOCUMENT
%
%
%
\begin{document}

%titlepage








\renewcommand{\authorDOC}{Daniel Sabalakov}

\renewcommand{\titleDOC}{Cracker Protein Lab}

\renewcommand{\classDOC}{DE Intro to Biology I}





\begin{titlepage}
  \titleDOC
  \classDOC
  \dateDOC
  \authorDOC
\end{titlepage}

\section{Abstract}\label{abstract}

To show the denaturing of amylase through the change in pH levels and how the efficacy and efficiency of enzyme performance are scale-based.

\section{Background / Introduction}\label{background-introduction}

\section{Question / Problem}\label{question-problem}

How does \textbf{\emph{amylase}}, an enzyme found in saliva, break polysaccharides into simpler carbohydrates in different PH levels?

\section{Elements \& Hypothesis}\label{elements-hypothesis}

\subsection{Hypothesis}\label{hypothesis}

If the pH levels of three solutions of pH \{4, 5.5, 7\}, are exposed to the enzyme \emph{amylase} found in saliva and mixed with a polysaccharide containing cracker, then there should be a trend whereas the pH 7 solution will have the highest density of glucose, the pH 5.5 solution should have the medium density of glucose, and the pH 4 solution should have the lowest density of glucose.

\subsection{Experimental Elements}\label{experimental-elements}

\begin{itemize}
\tightlist
\item
  Independent Variable: PH Level of the solution
\item
  Dependent Variable: Glucose concentration produced (in mmol/L)
\item
  Experimental Groups:

  \begin{itemize}
  \tightlist
  \item
    Group 1: Saltine Cracker, Saliva, pH 4.0 solution
  \item
    Group 2: Saltine Cracker, Saliva, pH 5.5 solution
  \end{itemize}
\item
  Positive Control Group:

  \begin{itemize}
  \tightlist
  \item
    Saltine Cracker, Saliva, optimal pH solution (pH 7.0)
  \end{itemize}
\item
  Negative Control Group:

  \begin{itemize}
  \tightlist
  \item
    Saltine Cracker, optimal pH solution (pH 7.0)
  \end{itemize}
\item
  Controls

  \begin{itemize}
  \tightlist
  \item
    Saliva to Solution ratio
  \item
    Cracker to Solution ratio
  \item
    Time allowed for enzymatic processes
  \item
    Method of glucose measurement
  \end{itemize}
\end{itemize}

\section{Materials / Set up}\label{materials-set-up}

\subsection{Materials}\label{materials}

The following materials were gathered:

\begin{enumerate}
\def\labelenumi{\arabic{enumi}.}
\tightlist
\item
  4 saltine cracker pieces (approximately 1 gram each)
\item
  Roughly 100 mL of Heinz all-natural vinegar
\item
  pH test strips
\item
  Roughly 250 mL of distilled water
\item
  8 15-25mL test tubes
\item
  Test tube rack
\item
  Life2O Pet Urinalysis Test Strips: Dog and Cat Urine Test Strips
\item
  Genie Vortex-Genie 2 Vortex Mixer SI-0236
\item
  Fresh Human Saliva (Roughly 20 mL total)
\item
  Stopwatch
\item
  Plastic pipettes
\item
  4 stirring rods
\item
  Plastic gloves
\item
  Tape
\item
  Black Marker
\item
  2x 250 mL beaker
\item
  Plastic test tubes
\item
  Mortar \& Pestile
\end{enumerate}

\subsection{Set Up Procedure}\label{set-up-procedure}

\begin{enumerate}
\def\labelenumi{\arabic{enumi}.}
\tightlist
\item
  All lab materials were collected and thoroughly inspected to ensure cleanliness.
\item
  Two sets of four test tubes were labeled with \{"C","pH=7","pH=5.5","pH=4"\} using tape and a marker.
\item
  The first set of four test tubes were filled with 10 mL of distilled water.

  \begin{itemize}
  \tightlist
  \item
    The water was filled up in the first beaker with about 100 mL of distilled water.
  \item
    The vinegar was filled up to about the 50 mL mark of the second beaker.
  \end{itemize}
\item
  Then, according to the test tube labels, the first set of four had 2 mL of vinegar added incrementally (until the pH was correct) and checked via a pH strip.
\item
  When all four test tubes of the first set were at the pH level indicated by each respective label, the tubes were put away, in a row, back onto the test tube rack
\end{enumerate}

An exemplar image of an example setup without pH solutions can be found below:

\begin{figure}[htbp]
	 \centering
	 \includegraphics[width=0.9\textwidth]{C:/Users/Scoresh/Desktop/Lab-Gen/Lab-Gen/lab-gen-source/images/bio-lab-temp.png}
\end{figure}

\section{Methods / Procedure}\label{methods-procedure}

\begin{enumerate}
\def\labelenumi{\arabic{enumi}.}
\tightlist
\item
  Saliva was collected by each individual labmate into their respective plastic test tube. This was to minimize the safety impact of having different contaminants on each test tube.
\item
\end{enumerate}

\section{Data / Results}\label{data-results}

\begin{longtable}[]{@{}cc@{}}
\toprule\noalign{}
Group & Glucose Measured (mmoL/L) \\
\midrule\noalign{}
\endhead
\bottomrule\noalign{}
\endlastfoot
Negative Control & 0 \\
Positive Control & 15 \\
pH=5.5 & 5 \\
pH=4.0 & 0 \\
\end{longtable}

\section{Conclusion / Discussion}\label{conclusion-discussion}



\end{document}
